\section{Study Design}
\label{sec:design}

\subsection{Research questions}
\label{sec:rqs}
\begin{description}
\item[RQ1] How does providing FEVER-designated gold evidence affect aggregate LLM fact-verification accuracy?
\item[RQ2] When evidence changes a previously correct prediction into an incorrect one, how frequent and consistent are those regressions across models?
\item[RQ3] What observable diagnostic failure modes characterize evidence-induced regressions under progressive disclosure?
\item[RQ4] Do those diagnostic conclusions survive a change in adjudicating model family?
\end{description}

\subsection{FEVER dataset and sampling}
We sample from the FEVER development set \citep{thorne-etal-2018-fever}.
The manuscript experiment is a balanced 10{,}000-claim extract: 5{,}000 Supported and 5{,}000 Refuted, seed 42, \texttt{NOT ENOUGH INFO} excluded.
Gold evidence is resolved from official FEVER Wikipedia shards into the text shown to the models.
Sampling and validation artifacts are committed with the frozen repository; third-party FEVER and Wikipedia text are not relicensed by the project MIT license.

\subsection{Models and paired evidence conditions}
Each claim is evaluated four times: GPT-5.4 and GPT-5.4-mini, each under \texttt{claim\_only} and \texttt{claim\_plus\_evidence}.
That yields 40{,}000 predictions.
The evidence condition uses the designated FEVER evidence, not an open-web retrieval step.
Original GPT runs stored labels only, so later analysis cannot recover confidence or token-level attributions.

\subsection{Transition definitions}
For each model, a claim has one of four paired transitions:
\emph{robust} (correct without and with evidence),
\emph{rescue} (wrong without, correct with),
\emph{resistant} (wrong in both),
or \emph{regression} (correct without, wrong with).
A unique-claim regression is a claim that regresses for at least one of the two GPT models.
Unique-claim counts are never treated as model-observations: 115 GPT-5.4 regressions and 151 mini regressions are 226 unique claims, not 266.

\subsection{Analysis v2}
Analysis v2 treats the 40{,}000-row experiment file as immutable source data.
It rebuilds the paired dataset, runs McNemar tests and bootstrap confidence intervals, extracts linguistic and evidence-structure features, embeds claim--evidence pairs, and applies two local NLI models (nli-deberta-v3-small and distilbert-base-uncased-mnli) as diagnostics, not as ground truth.
Leakage audits and an independent 69-check headline verifier are part of the frozen record.
Confirmatory tests were pre-specified; NLI and feature analyses are exploratory.

\subsection{Diagnostic cohort}
The adjudication cohort contains 1{,}061 unique claims: all 226 unique regressions, 335 resistant, 250 rescue, and 250 robust.
One item (\texttt{SA-000352}) is provider-filtered, so judged $N=1{,}060$.
The Claude full panel judges only the 226 regressions.
The Claude residual panel judges only the 231 Cursor/Grok Stage~C leftovers.
Those three samples answer different questions and are not interchangeable.

\subsection{Progressive-disclosure adjudication}
\label{sec:protocol}
Five isolated judges see nested evidence representations:
Stage~A, gold evidence sentences only;
Stage~B, Stage~A plus page titles;
Stage~C, reconstructed structured FEVER evidence (chosen set, alternatives, archive flags).
Judges do not browse, retrieve new pages, or see FEVER labels, GPT predictions, cohort tags, NLI scores, or other judges.
Stages are frozen before the next packets are built.
Consensus and the three-way taxonomy are computed by code.
We call the categories \emph{diagnostic failure modes}: they describe where a blinded panel's warrant changes, not a causal mechanism inside GPT.

The primary panel locks \texttt{cursor-grok-4.6-high-fast} for all judges and a same-family resolver on non-high-consensus items.
Taxonomy comparison with Claude uses consensus labels, which are computed identically across panels.

\subsection{Cross-family replication}
\label{sec:cross-family-design}
GLM-5.3 contributes a completed Stage~A panel compared to frozen Cursor Stage~A ($n=1{,}060$).
Partial GLM Stages~B and~C exist only as provenance and are excluded from A$\rightarrow$B$\rightarrow$C mathematics.

The Claude residual experiment is a Stage~C robustness test on the 231 Cursor/Grok Ambiguous/Unresolved leftovers (\texttt{claude-opus-5-thinking-high}, five judges, 1{,}155 judgments, no resolver).
It is selection-conditioned on Grok leftovers.

The main Claude experiment independently reruns A$\rightarrow$B$\rightarrow$C on all 226 unique regressions (226$\times$5$\times$3 $= 3{,}390$ judgments), with fresh opaque IDs and no splicing of residual-panel judgments.
During execution, the first Stage~A batch-01 wave used judge-facing paths that named the panel directory and therefore the cohort-selection criterion.
Those 385 judgments were quarantined unread, paths were replaced with a neutral \texttt{blind\_io} namespace, and the batch was rerun under the same locked model, rubric, and evidence packets.
A later usage-limit interruption paused Stage~B; work resumed on the same Claude slug, and already-valid judgments were retained.
Overlapping recovery produced three duplicate Stage~B judge-batches; a content-independent first-ingested-wins rule kept one copy.
Details are in Appendix~\ref{app:deviations}.

\subsection{Statistical analysis and verification}
Paired accuracy uses McNemar tests.
Regression and rescue rates use bootstrap confidence intervals.
Cross-family taxonomy agreement is reported as exact match, broad-category match, and Cohen's $\kappa$.
Stage-wise ambiguity differences use McNemar tests on the paired 226-item sample.
The frozen artifact verifies Analysis~v2 (69/69), Cursor/Grok (33/33), read-only GLM Stage~A, Claude residual (38/38), and Claude full (21/21).
Cohen's $\kappa=0.537$ is taken from the canonical synthesis, not from the Claude \texttt{HEADLINES.json} registry.
